% Generated draft: requires mathematical and source review.
\documentclass[11pt]{article}
\usepackage[margin=25mm]{geometry}
\usepackage{fontspec}
\setmainfont{Latin Modern Roman}
\usepackage{amsmath,amssymb,mathtools,unicode-math}
\setmathfont{Latin Modern Math}
\usepackage{enumitem,needspace,xurl,hyperref}
\setlength{\parindent}{0pt}
\setlength{\parskip}{5pt}
\setlength{\emergencystretch}{3em}
\newcommand{\N}{\mathbb{N}}
\newcommand{\Z}{\mathbb{Z}}
\newcommand{\Q}{\mathbb{Q}}
\newcommand{\R}{\mathbb{R}}
\newcommand{\C}{\mathbb{C}}
\newcommand{\F}{\mathbb{F}}
\newcommand{\A}{\mathbb{A}}
\newcommand{\PP}{\mathbb{P}}
\newcommand{\E}{\mathbb{E}}
\newcommand{\eps}{\varepsilon}
\newcommand{\dd}{\,\mathrm{d}}
\begin{document}
\section*{UnsolvedMath ID: 30003984}
\textbf{Source problem code:} catalogue entry
\subsection{Definitions and mathematical statement}
Let \(d,n\) be positive integers with \(n\ge d\). A complex equiangular tight frame of \(n\) unit vectors in \(\C^d\) is a family \((\varphi_i)_{i=1}^n\) such that, for some real number \(c\ge0\),
\[
\varphi_i^*\varphi_i=1\quad(1\le i\le n),\qquad
|\varphi_i^*\varphi_j|^2=c\quad(1\le i\ne j\le n),\qquad
\sum_{i=1}^n\varphi_i\varphi_i^*=\frac nd I_d.
\]
Here \(^*\) denotes conjugate transpose and \(I_d\) is the identity operator on \(\C^d\). Write \(\mathcal E_{\C}(d,n)\) for the set of such families. When \(n=1\), necessarily \(d=1\), and the pairwise equiangularity condition is vacuous. These definitions do not specify the parameter ranges of the conjectures named below.

\emph{Problem as supplied.} Develop finite-field constructions that resolve difficult existence questions for complex equiangular tight frames, including the questions identified in the record as the Fickus and Zauner conjectures, and determine whether the resulting finite-field solutions can be transferred to complex equiangular tight frames.

More precisely, a fully specified version must identify the target pairs \((d,n)\), any additional structure required by each conjecture, and a finite-field notion of equiangular tight frame. For each intended target, the requested programme is to obtain the corresponding finite-field solutions and establish a transfer result, under explicitly stated hypotheses, that yields an element of \(\mathcal E_{\C}(d,n)\) with any required additional structure.

\emph{Unresolved qualifications.} The supplied record does not state either named conjecture, its parameter range, or its possible symmetry requirements. It also does not specify the finite fields, bilinear or Hermitian forms, finite-field equiangularity and tightness constants, allowable characteristics, or the meaning and hypotheses of lifting. In particular, the complex normalization \(n/d\) cannot simply be read as a finite-field scalar when the characteristic divides \(d\). No unconditional implication from existence over an arbitrary finite field to complex existence is asserted here. The record therefore specifies a research programme rather than a uniquely determined quantified conjecture.
\subsection{Short English statement}
The problem asks whether finite-field constructions can settle specified complex equiangular tight frame existence questions, including the named Fickus and Zauner conjectures, and whether those constructions can be transferred to \(\C\). The exact conjectures, finite-field models, and transfer hypotheses are missing, so the statement remains a qualified research programme requiring mathematical review.
\subsection{Sources}
\begin{itemize}
\item Local UnsolvedMath registry: supplied complete record for permanent ID \(30003984\), title \emph{Finite-Field Constructions of Complex Equiangular Tight Frames}, source problem code \texttt{catalogue entry}, in the Oberwolfach Reports Open Problems collection. No individual proposer is supplied; the proposed year is recorded as 2018.
\item Supplied report citation: \emph{Mini-Workshop: Algebraic, Geometric, and Combinatorial Methods in Frame Theory} (2019), Oberwolfach Reports, pp.\ 2769--2803, DOI \texttt{10.4171/owr/2018/46}. \url{https://doi.org/10.4171/owr/2018/46}. The record also labels this reference \emph{Oberwolfach Report 46/2018}. No precise passage for this problem is supplied.
\item Supplied literature reference: \emph{Frames over Finite Fields: Basic Theory}. \url{https://arxiv.org/abs/2012.12977}.
\item Supplied literature reference: \emph{Frames over Finite Fields: Orthogonal Geometry}. \url{https://arxiv.org/abs/2012.13642}.
\item Literature assessment dated exactly \texttt{2026-08-21}: the supplied curation labels the problem \texttt{partially\_solved}. It reports substantial finite-field analogues and transfer from a complex equiangular tight frame to infinitely many finite fields, while describing reverse transfer and the general existence questions as unresolved. This is supplied dated curation derived from the OWR record, not an independently verified literature assessment; its time-sensitive status and cited claims require independent checking.
\item The registry labels the statement \texttt{corrected\_verified}, with the explanation that mathematical content was preserved while notation, typography, and sentence structure were repaired for publication. This supplied editorial label does not establish independent mathematical verification of this draft.
\end{itemize}
\end{document}
